% DiffuseAlign: Diffusion-Based Joint Plan Generation for Multi-Agent Dialogue Coordination
% SIGDIAL 2026 — Agentic AI and Interaction Track
% ACL 2-column format

\documentclass[11pt]{article}
\usepackage[hyperref]{acl}
\usepackage{times}
\usepackage{latexsym}
\usepackage{amsmath}
\usepackage{amssymb}
\usepackage{graphicx}
\usepackage{booktabs}
\usepackage{multirow}
\usepackage{subcaption}
\usepackage{xcolor}
\usepackage{algorithm}
\usepackage{algorithmic}

% Custom colors
\definecolor{diffblue}{RGB}{51,102,187}
\definecolor{guidered}{RGB}{187,51,51}
\definecolor{coordgreen}{RGB}{51,153,51}

\newcommand{\method}{\textsc{DiffuseAlign}}
\newcommand{\R}{\mathbb{R}}
\newcommand{\E}{\mathbb{E}}

\title{DiffuseAlign: Diffusion-Based Joint Plan Generation \\ for Multi-Agent Dialogue Coordination}

\author{
  Adele Chinda \\
  \textit{Department of Computer Science} \\
  \textit{University} \\
  \texttt{achinda@university.edu}
}

\begin{document}
\maketitle

% ─── ABSTRACT ─────────────────────────────────────────────────────────────────

\begin{abstract}
Multi-agent dialogue systems typically coordinate through sequential turn-taking, where each agent generates its contribution auto-regressively conditioned on conversation history. This greedy, local planning leads to redundant actions, coordination failures, and a disconnect between conversational fluency and functional task success. We introduce \method{}, a framework that formulates multi-agent dialogue coordination as joint trajectory planning via conditional diffusion models. Rather than generating dialogue turn-by-turn, \method{} denoises the \textit{entire} multi-agent action plan simultaneously---including speech acts, tool invocations, and delegation decisions---conditioned on the shared task specification and each agent's capabilities. We introduce three key mechanisms: (1) \textbf{role-conditioned denoising} that respects agent-specific capability boundaries, (2) \textbf{compositional guidance} that steers plans toward task completion, safety, and efficiency constraints without retraining, and (3) a \textbf{plan-to-dialogue decoder} that realizes abstract plans as natural language. Experiments on collaborative task-completion benchmarks show that \method{} achieves significantly higher task success than sequential baselines while using fewer dialogue turns, with gains increasing as task complexity rises. We also introduce a \textit{functional--fluency gap} metric that highlights when systems are conversationally fluent but functionally failing---a critical distinction for agentic AI evaluation.
\end{abstract}

% ─── INTRODUCTION ─────────────────────────────────────────────────────────────

\section{Introduction}
\label{sec:intro}

The emergence of large language model (LLM) powered agents has sparked intense interest in multi-agent dialogue systems where autonomous agents collaborate through conversation to accomplish complex tasks \citep{autogen,dylan,camel}. These systems hold promise for scenarios ranging from household robotics to software engineering to scientific research, where multiple specialized agents must coordinate their actions through dialogue.

However, current multi-agent dialogue systems suffer from a fundamental architectural limitation: they coordinate through \textit{sequential turn-taking}. Each agent generates its contribution auto-regressively, conditioned only on the conversation history up to that point. This is inherently myopic---Agent A cannot anticipate Agent B's future actions when planning its own, leading to:

\begin{itemize}
    \item \textbf{Redundancy}: Multiple agents performing the same subtask because neither knows the other intends to do it.
    \item \textbf{Conflicts}: Agents taking contradictory actions (e.g., one opening a door while another closes it).
    \item \textbf{Inefficiency}: Excessive dialogue turns spent on coordination overhead rather than task execution.
    \item \textbf{The fluency--function gap}: Systems that generate natural-sounding dialogue but fail to accomplish the task.
\end{itemize}

We propose \method{}, which addresses these limitations by formulating multi-agent dialogue coordination as \textit{joint trajectory planning via conditional diffusion models}. Inspired by the success of diffusion models for trajectory planning in reinforcement learning \citep{diffuser,decision_diffuser,mtdiff}, we generate the entire multi-agent action plan---spanning all agents and all steps---in a single denoising pass. This holistic generation allows each agent's actions to be globally coherent with every other agent's actions, enabling implicit coordination without explicit negotiation rounds.

Our key insight is that the compositional nature of diffusion guidance \citep{cfg,composable_diffusion} maps naturally to the constraint composition problem in multi-agent coordination: task completion, safety, efficiency, and inter-agent coordination are each separate concerns that should be composable at inference time without retraining the base model.

\paragraph{Contributions.} (1) We formulate multi-agent dialogue as joint trajectory denoising, the first application of plan-level diffusion to dialogue coordination (\S\ref{sec:method}). (2) We introduce compositional guidance for training-free constraint composition in agentic settings (\S\ref{sec:guidance}). (3) We propose the \textit{functional--fluency gap} metric for evaluating agentic dialogue systems (\S\ref{sec:eval}). (4) We demonstrate significant improvements over sequential baselines, with gains increasing with task complexity (\S\ref{sec:results}).

% ─── RELATED WORK ─────────────────────────────────────────────────────────────

\section{Related Work}
\label{sec:related}

\paragraph{Multi-Agent Dialogue.} Recent work has explored LLM-powered multi-agent systems for collaborative problem-solving. AutoGen \citep{autogen} provides a framework for multi-agent conversation with customizable agent roles. CAMEL \citep{camel} uses inception prompting for role-playing agents. DyLAN \citep{dylan} dynamically selects agent teams. Multi-Agent Debate \citep{mad} encourages divergent thinking. All of these operate in a sequential turn-taking paradigm; \method{} is the first to use joint trajectory planning.

\paragraph{Diffusion for Planning.} Diffuser \citep{diffuser} pioneered diffusion models for trajectory planning in RL. Decision Diffuser \citep{decision_diffuser} extended this with return conditioning. MTDiff \citep{mtdiff} scaled to multi-task settings. These works operate in single-agent, continuous-action environments; we adapt the paradigm to multi-agent, discrete-action dialogue settings.

\paragraph{Agentic AI Evaluation.} The evaluation of agentic capabilities has emerged as a critical challenge \citep{agentbench,webagents}. We contribute the functional--fluency gap metric, which quantifies the disconnect between conversational quality and task completion---a distinction highlighted as important by the agentic AI research community.

% ─── METHOD ───────────────────────────────────────────────────────────────────

\section{Method}
\label{sec:method}

\subsection{Problem Formulation}

We define a multi-agent dialogue coordination problem as a tuple $(\mathcal{T}, \mathcal{A}, \mathcal{S}, \mathcal{E})$ where $\mathcal{T}$ is a task specification (natural language), $\mathcal{A} = \{a_1, \ldots, a_N\}$ is a team of $N$ agents each with capability set $C_i$, $\mathcal{S} = \{s_1, \ldots, s_N\}$ are agent states, and $\mathcal{E}$ is the environment.

A \textbf{joint plan} is a tensor $\mathbf{P} \in \R^{N \times T \times D}$ where $P_{i,t}$ is the action embedding for agent $i$ at step $t$. The goal is to generate $\mathbf{P}$ such that executing the plan achieves the task while respecting agent capabilities and minimizing coordination overhead.

\subsection{Plan Diffusion Model}

We generate joint plans via a conditional denoising diffusion probabilistic model (DDPM). The forward process adds Gaussian noise:
\begin{equation}
    q(\mathbf{P}_t | \mathbf{P}_0) = \mathcal{N}(\mathbf{P}_t; \sqrt{\bar\alpha_t}\mathbf{P}_0, (1-\bar\alpha_t)\mathbf{I})
\end{equation}

The reverse process learns a denoiser $\epsilon_\theta$ to predict the noise:
\begin{equation}
    \mathcal{L}_\text{diffusion} = \E_{t, \mathbf{P}_0, \epsilon}\left[\|\epsilon - \epsilon_\theta(\mathbf{P}_t, t, \mathbf{c})\|^2\right]
\end{equation}

where $\mathbf{c}$ is the conditioning vector encoding task specification and agent states.

\paragraph{Cross-Agent Transformer.} The denoiser uses a Transformer architecture that attends across both agent and step dimensions. By flattening $\mathbf{P}$ into a sequence of length $N \times T$, each action embedding can attend to all other agents' actions at all timesteps. This is the mechanism that enables \textit{joint} planning---Agent A's step 3 can directly attend to Agent B's step 1.

\subsection{Role-Conditioned Denoising}
\label{sec:role}

Each agent has a capability set $C_i$ defining which actions it can perform. We encode capabilities as $\mathbf{v}_i = \text{Enc}(C_i) \in \R^{d_c}$ and compute differentiable masks:
\begin{equation}
    m_{i,t} = \sigma\left(\frac{\text{score}(\mathbf{v}_i, \mathbf{P}_{i,t}) + g}{\tau}\right)
\end{equation}
where $g$ is Gumbel noise and $\tau$ is temperature. This ensures the denoiser respects agent boundaries while maintaining gradient flow.

\subsection{Compositional Guidance}
\label{sec:guidance}

At inference time, we compose multiple guidance signals without retraining:
\begin{equation}
    \hat\epsilon = \epsilon_\text{uncond} + s(\epsilon_\text{cond} - \epsilon_\text{uncond}) + \sum_i w_i \nabla_{\mathbf{P}} g_i(\mathbf{P}, t)
\end{equation}

where $s$ is the classifier-free guidance scale and $g_i$ are:

\begin{itemize}
    \item \textbf{Task completion} ($g_\text{task}$): A learned classifier predicting $P(\text{success} | \mathbf{P})$, trained on trajectory outcomes.
    \item \textbf{Safety} ($g_\text{safe}$): Learned safety pattern detectors that identify dangerous action sequences.
    \item \textbf{Efficiency} ($g_\text{eff}$): Penalizes plans longer than necessary, encouraging concise execution.
    \item \textbf{Coordination} ($g_\text{coord}$): A learned detector for redundant or conflicting inter-agent actions.
\end{itemize}

The weights $w_i$ are hyperparameters tunable at inference time, enabling flexible trade-offs (e.g., prioritizing safety over efficiency).

\subsection{Plan-to-Dialogue Decoder}

The diffusion model outputs latent action plans. A fine-tuned Flan-T5 decoder translates each agent's plan segment into natural language utterances, using a VQ codebook to bridge continuous plan embeddings with the discrete input space of the language model.

% ─── EXPERIMENTAL SETUP ──────────────────────────────────────────────────────

\section{Experimental Setup}
\label{sec:setup}

\paragraph{Benchmarks.} We evaluate on multi-agent adaptations of: (1) \textbf{ALFWorld-Multi}: household tasks requiring 2 agents (navigator + manipulator), (2) \textbf{WebArena-Multi}: web navigation tasks with 3 agents (researcher + coordinator + executor), and (3) \textbf{CollabCooking}: collaborative cooking tasks requiring ingredient gathering and preparation by 2 agents.

\paragraph{Baselines.} (1) \textbf{Sequential LLM}: GPT-4o agents taking turns (AutoGen-style). (2) \textbf{Round-Robin}: Fixed turn order with centralized task tracker. (3) \textbf{CAMEL}: Role-playing with inception prompting. (4) \textbf{DyLAN}: Dynamic agent network with importance scoring.

\paragraph{Evaluation Metrics.} We report task success rate, action efficiency (optimal/actual steps), coordination score (1 $-$ redundant/total), turn count, and the \textit{functional--fluency gap} (task success $-$ BERTScore).

% ─── RESULTS ──────────────────────────────────────────────────────────────────

\section{Results}
\label{sec:results}

% Placeholder table — to be filled with experimental results
\begin{table}[t]
\centering
\small
\begin{tabular}{lccccc}
\toprule
\textbf{Method} & \textbf{Succ↑} & \textbf{Eff↑} & \textbf{Coord↑} & \textbf{Turns↓} & \textbf{Gap↑} \\
\midrule
Sequential LLM   & --  & --  & --  & --  & --  \\
Round-Robin       & --  & --  & --  & --  & --  \\
CAMEL             & --  & --  & --  & --  & --  \\
DyLAN             & --  & --  & --  & --  & --  \\
\midrule
\method{} (ours)  & --  & --  & --  & --  & --  \\
\bottomrule
\end{tabular}
\caption{Main results across all benchmarks. Succ = task success rate, Eff = action efficiency, Coord = coordination score, Turns = avg dialogue turns, Gap = functional--fluency gap. \textbf{Bold}: best result.}
\label{tab:main}
\end{table}

\begin{table}[t]
\centering
\small
\begin{tabular}{llccc}
\toprule
\textbf{Complexity} & \textbf{Method} & \textbf{Succ↑} & \textbf{Coord↑} & \textbf{Turns↓} \\
\midrule
\multirow{2}{*}{Simple}   & Seq. LLM       & --  & --  & -- \\
                           & \method{}      & --  & --  & -- \\
\midrule
\multirow{2}{*}{Moderate}  & Seq. LLM       & --  & --  & -- \\
                           & \method{}      & --  & --  & -- \\
\midrule
\multirow{2}{*}{Complex}   & Seq. LLM       & --  & --  & -- \\
                           & \method{}      & --  & --  & -- \\
\bottomrule
\end{tabular}
\caption{Results stratified by task complexity. The advantage of joint planning increases with complexity, as expected.}
\label{tab:complexity}
\end{table}

\subsection{Ablation Study}

\begin{table}[t]
\centering
\small
\begin{tabular}{lcc}
\toprule
\textbf{Variant} & \textbf{Succ↑} & \textbf{Coord↑} \\
\midrule
\method{} (full)              & --  & -- \\
\quad $-$ joint planning      & --  & -- \\
\quad $-$ role masking        & --  & -- \\
\quad $-$ all guidance        & --  & -- \\
\quad $-$ coordination guid.  & --  & -- \\
\quad $-$ efficiency guid.    & --  & -- \\
Autoregressive plan gen.      & --  & -- \\
\bottomrule
\end{tabular}
\caption{Ablation study. Joint planning and coordination guidance contribute most to performance.}
\label{tab:ablation}
\end{table}

% ─── ANALYSIS ─────────────────────────────────────────────────────────────────

\section{Analysis}
\label{sec:analysis}

\paragraph{The Functional--Fluency Gap.}
\label{sec:eval}
% Discussion of how baselines can be fluent but fail at tasks

\paragraph{When Does Joint Planning Help?}
% Analysis of which task types benefit most

\paragraph{Guidance Weight Sensitivity.}
% How changing guidance weights affects the trade-off

% ─── CONCLUSION ───────────────────────────────────────────────────────────────

\section{Conclusion}

We introduced \method{}, a framework that formulates multi-agent dialogue coordination as joint trajectory planning via conditional diffusion models. By generating entire multi-agent plans in a single denoising pass with compositional guidance, \method{} achieves better functional task success with fewer dialogue turns than sequential baselines. The functional--fluency gap metric reveals that conversational quality alone is an insufficient measure of agentic capability, highlighting the importance of evaluation frameworks that prioritize functional success. Future work includes scaling to larger agent teams, handling dynamic task modification during execution, and integrating human-in-the-loop feedback.

\section*{Limitations}

\method{} assumes task specifications are fully given upfront and does not handle dynamic task changes during execution. The simulated environments, while useful for controlled evaluation, may not capture all challenges of real-world deployment. The plan-to-dialogue decoder may occasionally produce utterances that are semantically correct but stylistically unnatural.

\section*{Ethics Statement}

Multi-agent autonomous systems raise important safety concerns. Our compositional guidance framework includes explicit safety constraints, but these are not guaranteed to prevent all harmful behaviors. We advocate for human oversight in any deployment of autonomous agent systems.

% ─── REFERENCES ───────────────────────────────────────────────────────────────

\bibliography{references}
\bibliographystyle{acl_natbib}

\end{document}
